\documentclass[10pt,a4paper]{amsart}
\usepackage[margin=19mm]{geometry}
\usepackage{amsmath,amssymb,mathtools}
\usepackage[scr=rsfs,cal=euler]{mathalfa}
\usepackage{xfrac}
\usepackage[unicode,hidelinks,bookmarks=false]{hyperref}
\newcommand{\ie}{i.e.\ }
\newcommand{\cf}{cf.\ }
\newcommand{\resp}{respectively\ }
\newcommand{\Subsection}{\subsection}
\providecommand{\nobreakdash}{\nobreak\hbox{-}}
\setlength{\emergencystretch}{2em}
\multlinegap0pt
\usepackage{amssymb}
\newtheorem{theorem}{Theorem}
\title{On some extensions of generalized counting processes}
\author{Lyudmyla Sakhno, Artem Storozhuk}
\date{Source: arXiv:2604.00882v1 (CC BY 4.0)}
\begin{document}
\maketitle

\noindent\emph{Source and licence.} On some extensions of generalized counting processes. Version arXiv:2604.00882v1 (CC BY 4.0), distributed by arXiv under the Creative Commons Attribution 4.0 International licence, http://creativecommons.org/licenses/by/4.0/, as recorded in the arXiv licence metadata for this exact version and shown on the abstract page https://arxiv.org/abs/2604.00882v1. Full licence notice: \texttt{licenses/arxiv-2604.00882-CC-BY-4.0.txt}.\par\smallskip

\noindent\textit{Editorial scope.} Counted unit: the article's theorem temp3 (source0.tex lines 808--818). The article's complete original proof of it is delivered with it (source0.tex lines 819--858, one \texttt{proof} environment of 40 lines). No mathematics is written, repaired or re-derived: the statement and the proof are the article's own lines, in the article's own notation, and no retained line is substituted or re-flowed.  Retained uncounted as the source-local context the proof needs: lines 134--136; lines 789--791.  Nothing was omitted from any retained span, so the package carries no omission record.  No retained line contains a citation, so the wrapper carries no bibliography and no bibliography entry.  No retained line contains a figure environment, an image-inclusion command or drawing code, so no image is delivered and the package declares no asset dependency.  Editorial formatting only, in addition: an AMS \texttt{amsart} preamble that restates the article's own declarations and macros used by the retained lines, namely the environments \texttt{theorem} on separate counters, together with the standard packages those lines need.  The retained environments print in this document as Theorem 1 in order of appearance, whereas the article prints the same items with the numbering of its own full document.  The complete native source of the article as distributed in the arXiv e-print for this exact version is bundled unchanged as \texttt{sources/197617/source0.tex}.
\begin{equation}\label{t-lapl-inv}
\mathcal{L}_t\left(\ell_f(t,x)\right)(r)=\frac{f(r)}{r}e^{-xf(r)}.
\end{equation}

\begin{equation}\label{exfDDt}
\mathbb{D}^{\alpha,\rho}_t u(t)=\frac{\alpha \rho^{\alpha}}{\Gamma(1-\alpha)}\frac{d}{dt}\int_{0}^{t}u(t-s)\Gamma(-\alpha,s)ds
\end{equation}

\begin{theorem}	
	{The density $l_{\alpha,\rho}(x,t)$  solves the time-fractional boundary-initial problem}
	\begin{equation}\label{temp3}
		\begin{cases}
			(\mathbb{D}^{2\alpha,\rho} + \mathbb{D}^{\alpha,\rho}) l_{\alpha, \rho}(x,t) = -\frac{\partial}{\partial x} l_{\alpha, \rho}(x,t), \quad x > 0,\,\, t > 0,\,\, 0 < \alpha < \frac{1}{2},  \\
			l_{\alpha, \rho}(x,0) = \delta(x), \\
			l_{\alpha, \rho}(0,t) = \frac{2\alpha \rho^{2\alpha}}{\Gamma(1-2\alpha)} \Gamma(-2\alpha, t) + \frac{\alpha \rho^\alpha}{\Gamma(1-\alpha)} \Gamma(-\alpha, t),
		\end{cases}
			\end{equation}
with the generalized R-L derivatives defined in \eqref{exfDDt}.
\end{theorem}	

	\begin{proof} Firsly, we find the double Laplace transform of the solution  to the problem \eqref{temp3}.
		Taking the $t$-Laplace transform of the equation \eqref{temp3} we have (we  will write below simply ${l}(x,s)$ omitting the subscript $_{\alpha,\rho}$):
		\begin{align*}
			f^{2\alpha, \rho}(s) \tilde{l}(x,s) + f^{\alpha, \rho}(s) \tilde{l}(x,s) = -\frac{\partial}{\partial x} \tilde{l}(x,s).
		\end{align*}
Then, with $x$-Laplace transform we obtain
		\begin{align*}
			(f^{2\alpha, \rho}(s) + f^{\alpha, \rho}(s)) \tilde{\tilde{l}}(\gamma, s) = \tilde{l}(0,s) - \gamma \tilde{\tilde{l}}(\gamma,s)
		\end{align*}
and, taking into account the boundary condition, we can calculate:
		\begin{align*}
			\tilde{l}(0,s) &= \int_0^{+\infty} e^{-st} l(0,t) dt = \int_0^{+\infty} e^{-st} (\nu_1(t) + \nu_2(t)) dt = \frac{f_1(s) + f_2(s)}{s},			
		\end{align*}
where we have denoted
$$\nu_1(t)=\frac{2\alpha \rho^{2\alpha}}{\Gamma(1-2\alpha)} \Gamma(-2\alpha, t), \,\, \nu_1(t)= \frac{\alpha \rho^\alpha}{\Gamma(1-\alpha)} \Gamma(-\alpha, t)$$	
and
$$f_1(s) = f^{2\alpha, \rho}(s) = (s+\rho)^{2\alpha} - \rho^{2\alpha}, \,\,
			f_2(s) = f^{\alpha, \rho}(s) = (s+p)^\alpha - \rho^\alpha.$$
Thus,				
		\begin{equation}\label{temp4}
			\tilde{\tilde{l}}(\gamma, s) = \frac{f_1(s) + f_2(s)}{s(f_1(s) + f_2(s) + \gamma)}
		\end{equation}
		
				On the other hand, let $l(t,x)$ be the density of the inverse process $\mathcal{L}^{\alpha, \rho}(t)$, then we can write
		\begin{equation}\label{(5)}
			l(t,x) = \frac{P(\mathcal{L}^{\alpha, \rho}(t) \in dx)}{dx} = -\frac{\partial}{\partial x} P \{ \mathcal{H}^{\alpha, \rho}(x) < t \} = -\frac{\partial}{\partial x} \int_0^t h(u,x) du,
		\end{equation}
where $h(t,x)$ is the probability density of the 	process	$\mathcal{H}^{\alpha, \rho}(t)=H_1^{2\alpha, \rho}(t) +  H_2^{\alpha,\rho}(t),$ which has the Laplace exponent $f_1(s) + f_2(s)$.

		In view of \eqref{(5)}, the double Laplace transform of $l(t,x)$ can be calculated as follows:
		\begin{align}
			\tilde{\tilde{l}}(\gamma, s) &= \int_0^{+\infty} e^{-\gamma x} \int_0^\infty e^{-st} \left[ -\frac{\partial}{\partial x} \int_0^t h(u,x) du \right] dt dx \nonumber\\
			&= - \int_0^\infty e^{-\gamma x} \frac{\partial}{\partial x} \int_0^\infty e^{-st} \int_0^t h(u,x) du dt dx = \nonumber\\
			&= -\frac{1}{s} \int_0^\infty e^{-\gamma x} \frac{\partial}{\partial x} \tilde{h}(s,x) dx = -\frac{1}{s} \int_0^\infty e^{-\gamma x} \frac{\partial}{\partial x} e^{-x(f_1(s) + f_2(s))} dx = \nonumber\\
			&= \frac{f_1(s) + f_2(s)}{s} \int_0^\infty e^{-\gamma x} e^{-x(f_1(s) + f_2(s))} dx = \frac{f_1(s) + f_2(s)}{s(f_1(s) + f_2(s) + \gamma)}.\label{LL}
		\end{align}
		
		Alternatively, we know  the expression \eqref{t-lapl-inv} for the  Laplace transform of the density of the inverse subordinator with respect to $t$ from which we can write the $t$-Laplace for the density  under consideration, and then applying $x$-transform we  get again \eqref{LL}.
		
		Therefore, the double Laplace transform of the density of the inverse subordinator $\mathcal{L}^{\alpha, p}(t)$ coincides with that of the solution of \eqref{temp3} given by \eqref{temp4}.\end{proof}

\end{document}
